\section{Multi-Agent-Entscheidung (Kap. 04)}
\label{sec:multi-agent}

StudyMummy ist in der untersuchten Implementierung als \emph{hierarchisches,
nachrichtengetriebenes Multi-Agent-System} einzuordnen. Entscheidend ist nicht
die Zahl der Prompts, sondern ausführbare Agentenmerkmale: Planning, Tutor und
Reviewer implementieren denselben \texttt{MASAgent}-Vertrag, besitzen je eine
Identität, ein eigenes Ziel, Capabilities und einen getrennten
\texttt{AgentLocalState}. Dieser Zustand zählt empfangene und gesendete
Nachrichten sowie Entscheidungen und hält eine kleine lokale Erinnerung für
den aktuellen Turn. Alle drei nehmen einen Nachrichtentyp wahr, entscheiden
innerhalb ihrer Rolle und erzeugen selbst Folgebotschaften.

Kooperation erfolgt über einen run-lokalen \texttt{MASMessageBus} mit
Performatives wie \texttt{request}, \texttt{delegate}, \texttt{propose},
\texttt{critique} und \texttt{accept}. Ein gemeinsames
\texttt{AgentBlackboard} repräsentiert die veränderliche Umgebung aus Plan,
Entwurf, Review und Toolbeobachtungen. Der normale Pfad lautet Planner--Tutor--Reviewer.
Bei Dissens routet der Reviewer autonom eine \texttt{revision\_request} an den
Tutor oder eine \texttt{replan\_request} an den Planner. Dadurch verändert er
den Kontrollfluss tatsächlich, statt lediglich einen Text zu ersetzen. Der
deterministische Orchestrator ist kein vierter Agent: Er dispatcht FIFO,
erfasst Laufzeiten und beendet nach höchstens der konfigurierten Rundenzahl.
Perception, Memory und RAG bleiben Umgebungsdienste, nicht künstlich
umbenannte Agenten.

\begin{center}
\begin{minipage}{\linewidth}
  \centering
  \begin{tikzpicture}[
    >=Latex,
    box/.style={draw,rounded corners,align=center,minimum height=8mm,text width=25mm,font=\small},
    role/.style={box,fill=hawblue!8},
    service/.style={box,fill=gray!10}
  ]
    \node[box] (ui) at (-4.8,1.7) {Lernende Person\\Angular};
    \node[box] (api) at (-1.6,1.7) {FastAPI\\\texttt{/agent/chat}};
    \node[service] (env) at (1.6,1.7) {Perception\\Session + RAG};
    \node[service] (memory) at (4.8,1.7) {Memory\\Chatlog + DB};
    \node[service,text width=31mm] (bus) at (0,0) {MAS Runtime\\Message Bus + Blackboard};
    \node[role] (planner) at (-4.2,-1.8) {Planner Agent\\Plan/Replan};
    \node[role] (tutor) at (0,-1.8) {Tutor Agent\\Act/Revise};
    \node[role] (reviewer) at (4.2,-1.8) {Reviewer Agent\\Accept/Critique};
    \node[service] (tools) at (-2.2,-3.6) {Tool Registry\\Side Effects};
    \node[service,text width=34mm] (llm) at (2.2,-3.6) {gemeinsamer OpenAI-kompatibler LLM-Adapter};

    \draw[->,semithick] (ui) -- (api);
    \draw[->,semithick] (api) -- (env);
    \draw[->,semithick] (env) -- (bus);
    \draw[->,semithick] (api) -- (memory);
    \draw[<->,semithick] (bus) -- (planner);
    \draw[<->,semithick] (bus) -- (tutor);
    \draw[<->,semithick] (bus) -- (reviewer);
    \draw[->,semithick] (tutor) -- (tools);
    \draw[dashed,->] (llm) -- (planner);
    \draw[dashed,->] (llm) -- (tutor);
    \draw[dashed,->] (llm) -- (reviewer);
  \end{tikzpicture}
  \captionsetup{hypcap=false}
  \captionof{figure}{Selbst erstellte Ist-Architektur des hierarchischen MAS.
    Rückpfeile am Bus stehen für Revision und Neuplanung. Volle
    Nachrichten-Payloads werden nicht an die API serialisiert.}
  \label{fig:mas-architektur}
\end{minipage}
\end{center}

\Cref{fig:mas-architektur} zeigt zugleich die Abgrenzung: Das MAS ist homogen,
seriell und in-process. Die Agenten teilen sich einen Modelladapter, besitzen
keine eigenen Prozesse und ihr lokaler Zustand endet mit dem Turn. Es gibt
weder Parallelität, freie Aushandlung noch horizontale Fehlerisolation. Ein
Single-Agent-Entwurf wäre günstiger, würde jedoch die getrennten Ziele,
Capability-Grenzen und den überprüfbaren Dissenspfad verlieren. Verteilte
Agentendienste wären umgekehrt ohne Parallelitäts- oder Verfügbarkeitsziel
unnötig komplex. Die gewählte Mitte liefert echte, aber bewusst begrenzte
Kooperation und bleibt für einen Tutor-Turn testbar.
